\documentclass[a4paper,11pt]{article}
        \usepackage[top=15mm,bottom=18mm,left=16mm,right=16mm]{geometry}
        \usepackage{fontspec}
        \usepackage{xeCJK}
        \setmainfont{Times New Roman}
        \setCJKmainfont{Yu Gothic}
        \setmonofont{Consolas}
        \setCJKmonofont{Yu Gothic}
        \usepackage{booktabs}
        \usepackage{longtable}
        \usepackage{tabularx}
        \usepackage{array}
        \usepackage{enumitem}
        \usepackage{hyperref}
        \usepackage{listings}
        \usepackage{xcolor}
        \usepackage{amsmath}
        \usepackage{amssymb}
        \hypersetup{
          colorlinks=true,
          linkcolor=blue,
          urlcolor=blue,
          pdftitle={preflight 判定ロジック},
          pdfauthor={Codex}
        }
        \lstset{
          basicstyle=\ttfamily\small,
          breaklines=true,
          columns=fullflexible,
          keepspaces=true,
          frame=single,
          framerule=0.2pt,
          rulecolor=\color{black},
          showstringspaces=false
        }
        \setlength{\parskip}{0.42em}
        \setlength{\parindent}{1em}
        \setlength{\emergencystretch}{3em}
        \renewcommand{\arraystretch}{1.15}
        \title{28. preflight 判定ロジック}
        \author{solar\_ws0129-main}
        \date{2026-07-29}
        \begin{document}
        \maketitle
        \tableofcontents
        \clearpage

        \section{対象}
        \begin{tabularx}{\linewidth}{>{\raggedright\arraybackslash}p{0.18\linewidth} X}
        \toprule
        項目 & 内容 \\
        \midrule
        ファイル & \path|mpc_solarcar/solar_preflight_logic.py| \\
        ソースSHA-256 & \path|2af5eb32f0e80a828a2fde431fe670d653b1aca06ee4c000690c8e503aa684d9| \\
        種別 & Python \\
        区分 & runtime helper \\
        役割 & 計測鮮度や command gate の純判定を Node 本体から切り出したロジック関数群。 \\
        起動文脈 & preflight と speed bridge の共通判定層。 \\
        \bottomrule
        \end{tabularx}

        \section{このファイルを読む理由}
        計測鮮度や command gate の純判定を Node 本体から切り出したロジック関数群。

        \begin{itemize}[leftmargin=1.6em]
        \item evaluate\_freshness と evaluate\_command\_gate が中心。
        \end{itemize}

        \section{呼び出し関係}
        \subsection{どこから呼ばれるか}
        \begin{itemize}[leftmargin=1.6em]
        \item \path|mpc_solarcar/solar_preflight_node.py|
\item \path|mpc_solarcar/speed_command_bridge_node.py|
        \end{itemize}

        \subsection{次に読むと理解が進むファイル}
        \begin{itemize}[leftmargin=1.6em]
        \item 特記事項なし。
        \end{itemize}

        \subsection{同一リポジトリ内の主要依存}
        \begin{itemize}[leftmargin=1.6em]
        \item 同一リポジトリ内の明示 import は少ない。
        \end{itemize}

        \section{ソース冒頭の把握}
        モジュール docstring または先頭意図の要約:

        モジュール docstring は明示されていない。

        \subsection{代表コード断片}
        \begin{lstlisting}[language=Python]
@dataclass(frozen=True)
class FreshnessResult:
    state: str
    health: float
    diagnostic: str


@dataclass(frozen=True)
class CommandGateResult:
    allowed: bool
    reason: str


def evaluate_freshness(
    *,
    elapsed_sec: float,
    ages_sec: dict[str, float | None],
    required: tuple[str, ...],
    timeout_sec: float,
    startup_grace_sec: float,
\end{lstlisting}


        \section{主要構造の読み取り}
        主要クラスは FreshnessResult, CommandGateResult。 主要関数は evaluate\_freshness, evaluate\_command\_gate。

        \subsection{主要クラス・関数}
            \begin{longtable}{p{0.07\linewidth} p{0.16\linewidth} p{0.25\linewidth} p{0.40\linewidth}}
            \toprule
            行 & 種別 & 名前 & 補足 \\
            \midrule
            7 & class & \path|FreshnessResult| &  \\
14 & class & \path|CommandGateResult| &  \\
19 & function & \path|evaluate_freshness| &  \\
43 & function & \path|evaluate_command_gate| &  \\
            \bottomrule
            \end{longtable}

        \subsection{ファイルを上から読んだときの定義順}
            \begin{longtable}{p{0.07\linewidth} p{0.84\linewidth}}
            \toprule
            行 & 内容 \\
            \midrule
            7 & クラス FreshnessResult を定義する。 \\
14 & クラス CommandGateResult を定義する。 \\
19 & 関数 evaluate\_freshness を定義する。 \\
43 & 関数 evaluate\_command\_gate を定義する。 \\
            \bottomrule
            \end{longtable}

        \subsection{import 群の詳細}
            \begin{longtable}{p{0.07\linewidth} p{0.33\linewidth} p{0.50\linewidth}}
            \toprule
            行 & import 文 & このファイルで読む理由 \\
            \midrule
            1 & \path|from __future__ import annotations| & 型ヒントの遅延評価を有効にし、自己参照型や forward reference を安全に扱うため。 このファイル内での使用位置は少ないか、間接利用である。 \\
3 & \path|from dataclasses import dataclass| & dataclasses から dataclass を読み込み、このファイルの処理を組み立てるため。 このファイル内での主な使用位置は L6, L13。 \\
            \bottomrule
            \end{longtable}

        \section{このファイルを読む前に必要な基礎知識}
        次の章は、構文やROS用語を既知と仮定しないための説明である。

        \subsection{プログラム、プロセス、メモリ、オブジェクトを区別する}
ソースファイルはディスク上の文字列であり、それ自体は走っていない。Python実行可能プログラムがソースを読み、OSがその実行を一つのプロセスとして管理する。プロセスには仮想メモリ、開いているファイル、スレッド、終了コードなどが対応する。


\begin{lstlisting}
ソースファイル -> Pythonインタプリタが読み込む -> OS上のプロセス -> Pythonオブジェクトをメモリに生成 -> 関数やcallbackを実行
\end{lstlisting}


「メモリ上に生成する」とは、実行中プロセスが使う記憶領域に、その値の型、属性、参照関係を表す実体を用意することである。変数は実体そのものというより、そのオブジェクトを指す名前として理解するとPythonの代入が読みやすい。


\begin{lstlisting}[language=python]
node = MPCNode()
alias = node
# nodeとaliasは同じオブジェクトを参照する。
\end{lstlisting}


一つのプロセスには複数スレッドを持たせられる。スレッドは同じプロセスのメモリを共有するため、受信callbackと最適化callbackが同じself.zを書き換える場合は実行順と排他を検討する必要がある。

\paragraph{根拠資料}
\begin{itemize}[leftmargin=1.6em]
\item \href{https://docs.python.org/3/tutorial/classes.html}{Python公式チュートリアル: Classes}
\end{itemize}

\subsection{名前、代入、型、参照}
Pythonの代入文は、右辺を先に評価し、その結果のオブジェクトへ左辺の名前を結び付ける。`x = f()`では、まずfを呼び、その戻り値が得られてからxが更新される。


`float(...)`、`int(...)`、`bool(...)`は型名を呼び出して値を変換する。外部CSV、YAML、ROS parameterから来た値は型が期待どおりとは限らないため、このリポジトリでは明示変換が多い。


`None`は値が無いことを示す単一のオブジェクト、`math.nan`は浮動小数点値ではあるが有効な数値ではないことを示す。両者は用途が異なるため、`is None`と`math.isfinite`を使い分ける。

\paragraph{根拠資料}
\begin{itemize}[leftmargin=1.6em]
\item \href{https://docs.python.org/3/tutorial/classes.html}{Python公式チュートリアル: Classes}
\end{itemize}

\subsection{関数、引数、戻り値、スコープ}
`def`は関数本体をその場で実行する文ではない。関数オブジェクトを作り、その名前を現在の名前空間へ登録する。`f`は関数そのもの、`f()`はその関数を呼んだ結果である。


仮引数は関数定義側の受け取り名、実引数は呼出側から渡す値である。位置引数、キーワード引数、既定値付き引数、`*args`、`**kwargs`は渡し方が異なる。


\begin{lstlisting}[language=python]
def clip_speed(value, minimum=0.0, maximum=110.0):
    return min(maximum, max(minimum, value))

v = clip_speed(120.0, maximum=90.0)
\end{lstlisting}


関数内で作った通常の名前はローカルスコープに属する。メソッドが`self.z`を更新すればオブジェクトの状態が残るが、単に`z`を更新しただけならその呼出しのローカル値である。

\paragraph{根拠資料}
\begin{itemize}[leftmargin=1.6em]
\item \href{https://docs.python.org/3/tutorial/classes.html}{Python公式チュートリアル: Classes}
\end{itemize}

\subsection{module、package、import、console\_scripts}
Pythonファイルはmoduleとして読み込める。複数moduleをディレクトリにまとめたものがpackageである。`from .model import SolarCarModel`の先頭の点は、現在と同じpackage内のmodel moduleを指す相対importである。


import時にはファイルのトップレベル文が上から一度実行される。`def`や`class`は関数・クラスを登録するが、その本体の通常処理は呼び出すまで走らない。


setuptoolsの`console\_scripts`は、端末で使う実行可能名と`package.module:function`を対応付けるインストール時メタデータである。生成される実行用ラッパーはmoduleをimportし、指定関数を呼び、戻り値を終了コードとして扱う小さな入口である。


\begin{lstlisting}
ROS launchのexecutable名 -> install済みconsole script -> mpc_solarcar.mpc_nodeをimport -> main()を呼ぶ
\end{lstlisting}

\paragraph{根拠資料}
\begin{itemize}[leftmargin=1.6em]
\item \href{https://setuptools.pypa.io/en/latest/userguide/entry_point.html}{setuptools公式: Entry Points / Console Scripts}
\item \href{https://docs.python.org/3/tutorial/classes.html}{Python公式チュートリアル: Classes}
\end{itemize}

\subsection{例外、try/finally、with、資源解放}
例外は通常の戻り値とは別経路で異常を呼出元へ伝える。`try/except`は想定した異常を処理し、`finally`は成功・失敗にかかわらず後始末を行う。


`with open(...) as f:`はcontext managerを使い、ブロックを出るとファイルを閉じる。CSVやログの破損を避けるため、開いた資源の所有者と閉じる場所を明確にする。


`except Exception: pass`はノードを止めない利点がある一方、入力異常を隠して原因追跡を難しくする。安全に関係する値では、少なくとも頻度制限付き警告、異常カウンタ、fallback状態のpublishを検討する。



\subsection{class、独自クラス、インスタンス、self、継承}
クラスは、保持するデータと、そのデータを扱う関数を一つの型としてまとめる設計単位である。「独自クラス」とはPythonやROS 2が最初から用意した型ではなく、このプロジェクトが目的に合わせて定義した新しい型を指す。


`class MPCNode(Node)`は、rclpyのNodeを基底クラスとして継承し、Nodeが持つpublisher、subscription、timer、parameterなどの機能にMPC固有機能を追加する。丸括弧内は関数引数ではなく基底クラス指定である。


`MPCNode()`はクラスオブジェクトを呼び出して新しいインスタンスを作る。Pythonは新しい実体を用意し、その実体を第1引数selfとして`\_\_init\_\_`へ渡す。`self`は予約語ではなく慣習的な名前だが、この慣習を守ることでコードの意味が共有される。


\begin{lstlisting}[language=python]
class Counter:
    def __init__(self):
        self.value = 0

    def increment(self):
        self.value += 1

counter = Counter()
counter.increment()
\end{lstlisting}


`super().\_\_init\_\_(...)`は継承元の初期化処理を呼ぶ。MPCNodeではこれを省くとROSノードとして必要な内部実体が作られず、`create\_publisher`などを正しく使えない。

\paragraph{根拠資料}
\begin{itemize}[leftmargin=1.6em]
\item \href{https://docs.python.org/3/tutorial/classes.html}{Python公式チュートリアル: Classes}
\item \href{https://docs.ros.org/en/humble/Tutorials/Beginner-CLI-Tools/Understanding-ROS2-Nodes/Understanding-ROS2-Nodes.html}{ROS 2 Humble公式: Understanding nodes}
\end{itemize}

\subsection{先頭アンダースコア、\_\_init\_\_、名前付けの作法}
`self.\_step\_solar`の先頭1個のアンダースコアは「クラス内部で使う実装詳細」という慣習であり、アクセスを強制禁止する機構ではない。外部コードから呼べるが、公開APIとして依存しない意思表示である。


`\_\_init\_\_`の前後2個のアンダースコアはPythonが定めた特殊メソッド名である。クラス生成時に自動で呼ばれる。任意の名前に前後2個を付けて独自仕様を作ることは避ける。


`\_`で始まるローカル変数は、未使用または外部へ見せない意図を示す場合がある。例えば`\_base\_csv`は戻り値として受け取る必要はあるが、その後の処理では使わないことを読者へ示す。

\paragraph{根拠資料}
\begin{itemize}[leftmargin=1.6em]
\item \href{https://docs.python.org/3/tutorial/classes.html}{Python公式チュートリアル: Classes}
\end{itemize}

\subsection{デコレータと@dataclass}
`@名前`は、直後に定義した関数またはクラスを別の関数へ渡し、その結果で定義名を置き換えるデコレータ構文である。`@Class`という一般構文があるのではなく、`@dataclass`など実際のデコレータ名を書く。


`@dataclass`は型注釈付きフィールドから`\_\_init\_\_`、`\_\_repr\_\_`、比較処理などを自動生成する。物理パラメータや解析結果のように、名前付きデータを一まとまりで運ぶ用途に向く。


\begin{lstlisting}[language=python]
from dataclasses import dataclass

@dataclass
class State:
    soc: float
    temperature_c: float

state = State(soc=0.8, temperature_c=30.0)
\end{lstlisting}


自動生成は検証を自動で保証しない。単位、許容範囲、相互依存制約は`\_\_post\_init\_\_`や別の検証関数で確認する必要がある。

\paragraph{根拠資料}
\begin{itemize}[leftmargin=1.6em]
\item \href{https://docs.python.org/3/library/dataclasses.html}{Python公式ライブラリ: dataclasses}
\item \href{https://docs.python.org/3/tutorial/classes.html}{Python公式チュートリアル: Classes}
\end{itemize}

\subsection{制御、状態、入力、モデル予測制御MPC}
制御対象の内部を表す状態をx、操作入力をu、外乱・予報をwとすると、離散モデルは`x[k+1] = f(x[k], u[k], w[k])`と書ける。ソーラーカーではSoC、電池温度、距離、速度などが状態候補、速度目標や駆動トルクが入力候補になる。


\[
\min_{u_0,\ldots,u_{N-1}} \sum_{k=0}^{N-1}\ell(x_k,u_k,w_k)+V_f(x_N)
\]

MPCは現在状態からNステップ先まで予測し、目的関数と制約を満たす入力系列を求める。ただし実際に適用するのは通常先頭入力だけで、次回は新しい実測状態から再び解く。これがreceding horizonである。


予測モデル、目的関数、制約、ホライズン、solver、初期値のどれかが変わると答えも変わる。「MPCを使う」だけでは仕様は決まらず、これらを単位付きで追う必要がある。

\paragraph{根拠資料}
\begin{itemize}[leftmargin=1.6em]
\item \href{https://docs.scipy.org/doc/scipy/reference/generated/scipy.optimize.minimize.html}{SciPy公式: scipy.optimize.minimize}
\end{itemize}



        \section{関数・クラスを上から順に解説}
        \subsection{L7 クラス FreshnessResult}
            \begin{itemize}[leftmargin=1.6em]
              \item 行範囲: L7--L10
              \item 所属: module直下
              \item 定義: \path|FreshnessResult(bases=none)|
              \item docstring: 明示されていない。
              \item このブロックが直接呼ぶ主な関数・メソッド: 特に明示されていない。
              \item 戻り値の要点: 戻り値が無いか、return 文が明示されていない。
              \item 主なローカル代入名: 特になし。
              \item 読み取る主なインスタンス属性: 特になし。
              \item 更新する主なインスタンス属性: 特になし。
              \item 明示的に送出する例外: 明示的raiseはない。
              \item 制御構造の規模: 条件分岐 0、ループ 0、try 0
            \end{itemize}
            \paragraph{この定義を読むためのPython構文}
            \begin{itemize}[leftmargin=1.6em]
            \item class文は新しい型を定義する。丸括弧内は継承する基底クラスである。
\item @で始まる行は、定義したクラスを加工するdecoratorである。
            \end{itemize}
            \paragraph{上から順の処理}
            \begin{enumerate}[leftmargin=1.8em]
            \item state に  を代入する。
\item health に  を代入する。
\item diagnostic に  を代入する。
            \end{enumerate}
            \paragraph{代表コード断片}
            \begin{lstlisting}[language=Python]
class FreshnessResult:
    state: str
    health: float
    diagnostic: str
\end{lstlisting}

\subsection{L14 クラス CommandGateResult}
            \begin{itemize}[leftmargin=1.6em]
              \item 行範囲: L14--L16
              \item 所属: module直下
              \item 定義: \path|CommandGateResult(bases=none)|
              \item docstring: 明示されていない。
              \item このブロックが直接呼ぶ主な関数・メソッド: 特に明示されていない。
              \item 戻り値の要点: 戻り値が無いか、return 文が明示されていない。
              \item 主なローカル代入名: 特になし。
              \item 読み取る主なインスタンス属性: 特になし。
              \item 更新する主なインスタンス属性: 特になし。
              \item 明示的に送出する例外: 明示的raiseはない。
              \item 制御構造の規模: 条件分岐 0、ループ 0、try 0
            \end{itemize}
            \paragraph{この定義を読むためのPython構文}
            \begin{itemize}[leftmargin=1.6em]
            \item class文は新しい型を定義する。丸括弧内は継承する基底クラスである。
\item @で始まる行は、定義したクラスを加工するdecoratorである。
            \end{itemize}
            \paragraph{上から順の処理}
            \begin{enumerate}[leftmargin=1.8em]
            \item allowed に  を代入する。
\item reason に  を代入する。
            \end{enumerate}
            \paragraph{代表コード断片}
            \begin{lstlisting}[language=Python]
class CommandGateResult:
    allowed: bool
    reason: str
\end{lstlisting}

\subsection{L19 関数 evaluate\_freshness}
            \begin{itemize}[leftmargin=1.6em]
              \item 行範囲: L19--L40
              \item 所属: module直下
              \item 定義: \path!evaluate_freshness(*, elapsed_sec: float, ages_sec: dict[str, float | None], required: tuple[str, ...], timeout_sec: float, startup_grace_sec: float) -> FreshnessResult!
              \item docstring: 明示されていない。
              \item このブロックが直接呼ぶ主な関数・メソッド: FreshnessResult, float, get, join
              \item 戻り値の要点: FreshnessResult('RUNNING', 1.0, 'solar telemetry and planner inputs are fresh') / FreshnessResult('STARTING', 0.25, 'waiting for required solar telemetry') / FreshnessResult('DEGRADED', 0.2, 'missing: ' + ', '.join(missing)) / FreshnessResult('DEGRADED', 0.4, 'stale: ' + ', '.join(stale))
              \item 主なローカル代入名: missing, name, stale
              \item 読み取る主なインスタンス属性: 特になし。
              \item 更新する主なインスタンス属性: 特になし。
              \item 明示的に送出する例外: 明示的raiseはない。
              \item 制御構造の規模: 条件分岐 3、ループ 0、try 0
            \end{itemize}
            \paragraph{この定義を読むためのPython構文}
            \begin{itemize}[leftmargin=1.6em]
            \item `->`以後は戻り値の型注釈であり、実行時の値を自動変換する命令ではない。
\item 内包表記は、反復・条件・値生成を一つの式で記述してcollectionを作る。
            \end{itemize}
            \paragraph{上から順の処理}
            \begin{enumerate}[leftmargin=1.8em]
            \item 条件 elapsed\_sec < startup\_grace\_sec を判定し、真なら内部処理を行う。
\item   FreshnessResult('STARTING', 0.25, 'waiting for required solar telemetry') を返す。
\item missing に [name for name in required if ages\_sec.get(name) is None] の結果を代入する。
\item stale に [name for name in required if ages\_sec.get(name) is not None and float(ages\_sec[name]) > timeout\_sec] の結果を代入する。
\item 条件 missing を判定し、真なら内部処理を行う。
\item   FreshnessResult('DEGRADED', 0.2, 'missing: ' + ', '.join(missing)) を返す。
\item 条件 stale を判定し、真なら内部処理を行う。
\item   FreshnessResult('DEGRADED', 0.4, 'stale: ' + ', '.join(stale)) を返す。
\item FreshnessResult('RUNNING', 1.0, 'solar telemetry and planner inputs are fresh') を返す。
            \end{enumerate}
            \paragraph{代表コード断片}
            \begin{lstlisting}[language=Python]
def evaluate_freshness(
    *,
    elapsed_sec: float,
    ages_sec: dict[str, float | None],
    required: tuple[str, ...],
    timeout_sec: float,
    startup_grace_sec: float,
) -> FreshnessResult:
    if elapsed_sec < startup_grace_sec:
        return FreshnessResult("STARTING", 0.25, "waiting for required solar telemetry")

    missing = [name for name in required if ages_sec.get(name) is None]
    stale = [
        name
        for name in required
        if ages_sec.get(name) is not None and float(ages_sec[name]) > timeout_sec
    ]
    if missing:
        return FreshnessResult("DEGRADED", 0.2, "missing: " + ", ".join(missing))
    if stale:
        return FreshnessResult("DEGRADED", 0.4, "stale: " + ", ".join(stale))
    return FreshnessResult("RUNNING", 1.0, "solar telemetry and planner inputs are fresh")
\end{lstlisting}

\subsection{L43 関数 evaluate\_command\_gate}
            \begin{itemize}[leftmargin=1.6em]
              \item 行範囲: L43--L68
              \item 所属: module直下
              \item 定義: \path!evaluate_command_gate(*, elapsed_sec: float, speed_input_age_sec: float | None, system_state: str, system_state_age_sec: float | None, startup_hold_sec: float, input_timeout_sec: float, system_state_timeout_sec: float, require_system_running: bool) -> CommandGateResult!
              \item docstring: 明示されていない。
              \item このブロックが直接呼ぶ主な関数・メソッド: CommandGateResult, max, str, strip, upper
              \item 戻り値の要点: CommandGateResult(True, 'ok') / CommandGateResult(False, 'startup\_hold') / CommandGateResult(False, 'missing\_speed\_command') / CommandGateResult(False, 'stale\_speed\_command')
              \item 主なローカル代入名: 特になし。
              \item 読み取る主なインスタンス属性: 特になし。
              \item 更新する主なインスタンス属性: 特になし。
              \item 明示的に送出する例外: 明示的raiseはない。
              \item 制御構造の規模: 条件分岐 7、ループ 0、try 0
            \end{itemize}
            \paragraph{この定義を読むためのPython構文}
            \begin{itemize}[leftmargin=1.6em]
            \item `->`以後は戻り値の型注釈であり、実行時の値を自動変換する命令ではない。
            \end{itemize}
            \paragraph{上から順の処理}
            \begin{enumerate}[leftmargin=1.8em]
            \item 条件 elapsed\_sec < max(0.0, startup\_hold\_sec) を判定し、真なら内部処理を行う。
\item   CommandGateResult(False, 'startup\_hold') を返す。
\item 条件 speed\_input\_age\_sec is None を判定し、真なら内部処理を行う。
\item   CommandGateResult(False, 'missing\_speed\_command') を返す。
\item 条件 speed\_input\_age\_sec > max(0.0, input\_timeout\_sec) を判定し、真なら内部処理を行う。
\item   CommandGateResult(False, 'stale\_speed\_command') を返す。
\item 条件 not require\_system\_running を判定し、真なら内部処理を行う。
\item   CommandGateResult(True, 'ok\_without\_system\_gate') を返す。
\item 条件 system\_state\_age\_sec is None を判定し、真なら内部処理を行う。
\item   CommandGateResult(False, 'missing\_system\_state') を返す。
\item 条件 system\_state\_age\_sec > max(0.0, system\_state\_timeout\_sec) を判定し、真なら内部処理を行う。
\item   CommandGateResult(False, 'stale\_system\_state') を返す。
\item 条件 str(system\_state).strip().upper() != 'RUNNING' を判定し、真なら内部処理を行う。
\item   CommandGateResult(False, 'system\_not\_running') を返す。
\item CommandGateResult(True, 'ok') を返す。
            \end{enumerate}
            \paragraph{代表コード断片}
            \begin{lstlisting}[language=Python]
def evaluate_command_gate(
    *,
    elapsed_sec: float,
    speed_input_age_sec: float | None,
    system_state: str,
    system_state_age_sec: float | None,
    startup_hold_sec: float,
    input_timeout_sec: float,
    system_state_timeout_sec: float,
    require_system_running: bool,
) -> CommandGateResult:
    if elapsed_sec < max(0.0, startup_hold_sec):
        return CommandGateResult(False, "startup_hold")
    if speed_input_age_sec is None:
        return CommandGateResult(False, "missing_speed_command")
    if speed_input_age_sec > max(0.0, input_timeout_sec):
        return CommandGateResult(False, "stale_speed_command")
    if not require_system_running:
        return CommandGateResult(True, "ok_without_system_gate")
    if system_state_age_sec is None:
        return CommandGateResult(False, "missing_system_state")
    if system_state_age_sec > max(0.0, system_state_timeout_sec):
        return CommandGateResult(False, "stale_system_state")
    if str(system_state).strip().upper() != "RUNNING":
        return CommandGateResult(False, "system_not_running")
    return CommandGateResult(True, "ok")
\end{lstlisting}











        \section{処理の流れ}
        \begin{enumerate}[leftmargin=1.7em]
        \item ソース中の主要関数を通じて処理を進める。
        \end{enumerate}

        \section{このファイルの位置づけ}
        この資料群では、\path|mpc_solarcar/solar_preflight_logic.py| を
        \textbf{runtime helper}の中核として扱う。
        したがって、ここで扱う処理は単独で完結せず、
        前段の profile / launch / telemetry と後段の planner / logger / report へ繋がる。

        \end{document}
